\documentclass[11pt]{article}
\usepackage[margin=1in]{geometry}
\usepackage{amsmath,amssymb,amsthm,hyperref,verbatim}
\newtheorem{theorem}{Theorem}
\newtheorem{lemma}{Lemma}
\newcommand{\E}{\mathbb E}
\title{A positive moment measure for the\\
Fennessey--Larcombe--French determinant sequence}
\author{Principia Math In Progress}
\date{11 September 2026}
\begin{document}
\maketitle
\begin{abstract}
We prove Conjecture 5.1 of Sun and Zhao: the sequence
$V_n^2-V_{n-1}V_{n+1}$, $n\geq2$, is infinitely log-convex.
In fact it is a Stieltjes moment sequence with an explicit positive
integral density after normalization. Positivity is proved analytically
on two endpoint intervals and by four rigorous, finite rational
certificates on the remaining interval. The certificate uses no
floating-point arithmetic.
\end{abstract}

\section*{Task 1: Resolve MathDB problem 332482}
Let $V_0=1$, $V_1=8$, and, for $n\geq2$, let
\begin{equation}\label{eq:Vrec}
 (n-1)n^2V_n
 =8(n-1)(3n^2-n-1)V_{n-1}-128(n-2)n^2V_{n-2}.
\end{equation}
Set $D_n=V_n^2-V_{n-1}V_{n+1}$ for $n\geq2$.
For a sequence $(a_m)_{m\geq0}$ define
$(\mathcal La)_m=a_ma_{m+2}-a_{m+1}^2$.
Infinite log-convexity means that every iterate $\mathcal L^r a$
is nonnegative. This is the usual convention with its initial index
shifted at each iteration.

\begin{theorem}\label{thm:main}
The sequence $(D_{m+2})_{m\geq0}$ is a Stieltjes moment sequence
represented by a finite positive measure on $[0,256]$ with support
equal to that interval. It is strictly infinitely log-convex.
In particular, published Conjecture 5.1 \cite[p.~601]{source} holds.
\end{theorem}

The MathDB title says Catalan--Larcombe--French, but its statement and
the published conjecture both specify the Fennessey--Larcombe--French
numbers $V_n$ used here.

\subsection*{A preliminary density identity}
We give the needed representation explicitly, including sequence
identification. Write
\[
 L=\frac2\pi,\quad M=2L-1,\quad
 c_j=\left(\frac{\binom{2j}j}{4^j}\right)^2,\quad
 K(m)=\int_0^{\pi/2}(1-m\sin^2\theta)^{-1/2}\,d\theta.
\]
Thus $K$ uses the parameter convention, not the modulus convention.
For $u\ne1/2$ in $(0,1)$ set
\begin{equation}\label{eq:rho}
 w(u)=\frac4{\pi^2}K(4u(1-u)),\qquad
 \rho(u)=w(u)+(2u-1)w'(u).
\end{equation}
Values at isolated singularities are irrelevant to all measures below.

\begin{lemma}\label{lem:rho}
The density $\rho$ is positive almost everywhere, symmetric about
$1/2$, and has mass $M<L$. For $n\geq1$,
\begin{equation}\label{eq:Vmoment}
 v_n:=\frac{V_n}{16^n}=L-\int_0^1u^n\rho(u)\,du.
\end{equation}
It has the positive series
\begin{equation}\label{eq:series}
 \rho(u)=\frac2\pi\sum_{j\geq1}\frac{j}{j+1}c_j
                 [4u(1-u)]^j.
\end{equation}
\end{lemma}

\begin{proof}
Let $X,Y$ be independent with density $1/(\pi\sqrt{1-x^2})$ on
$(-1,1)$, and put $U=(1+XY)/2$. The product-density substitution gives
\[
 f_{XY}(t)=\frac2{\pi^2}\int_{|t|}^1
 \frac{dx}{\sqrt{(1-x^2)(x^2-t^2)}}
 =\frac2{\pi^2}K(1-t^2).
\]
The second equality follows by setting
$x^2=t^2+(1-t^2)\sin^2\theta$ and reversing the angle.
Consequently $w$ is the probability density of $U$ and
\[
 p_n:=\E(U^n)=2^{-n}\sum_{j=0}^{\lfloor n/2\rfloor}\binom n{2j}c_j.
\]
Here the arcsine even moments are $\binom{2j}j/4^j$.
To verify the recurrence identification, let $H(q)=\sum c_jq^j$.
The coefficient ratio gives
$q(1-q)H''+(1-2q)H'-H/4=0$.
The generating function of $p_n$ is
$A(z)=2(2-z)^{-1}H(z^2/(2-z)^2)$, whence, by substitution,
\[
 z(1-z)(2-z)A''+(2-6z+3z^2)A'+(z-1)A=0.
\]
Thus $P_n=16^np_n$ satisfies, for $n\geq2$,
\[
 E_n:=n^2P_n-8(3n^2-3n+1)P_{n-1}+128(n-1)^2P_{n-2}=0,
\]
with $P_0=1,P_1=8$ and $E_1=P_1-8P_0=0$.
Set $W_n=(2n+1)P_n-16nP_{n-1}$ for $n\geq1$.
The residual of \eqref{eq:Vrec} with $W$ in place of $V$ is
$(n-1)(2n+1)E_n-16n^2E_{n-1}$.
It vanishes, and $W_1=8$, so $W_n=V_n$. In particular,
\begin{equation}\label{eq:weighted}
 v_n=\int_0^1\bigl[u^n+n(2u-1)u^{n-1}\bigr]w(u)\,du.
\end{equation}

Expanding $K(m)=(\pi/2)\sum c_jm^j$ in \eqref{eq:rho} gives
\eqref{eq:series}; indeed
$\rho=(4/\pi^2)[K(m)-4(1-m)K'(m)]$ with $m=4u(1-u)$.
The series proves positivity and symmetry.
The elementary bound $c_j\leq1/(j+1)$, obtained by induction from
the coefficient ratio, gives $K(m)=O(1+|\log(1-m)|)$.
Hence $h(u)=(2u-1)w(u)$ extends continuously through $1/2$ with
value zero. Its endpoint values are $h(0)=-L,h(1)=L$, and away
from $1/2$ its derivative is $h'=w+\rho>0$.
Integrating on each side and taking limits gives
$\int_0^1(w+\rho)=2L$. Thus $\rho$ is integrable of mass
$M=2L-1$, and $h$ is absolutely continuous.
Integration by parts in \eqref{eq:weighted} now gives
\eqref{eq:Vmoment}. The boundary contribution at zero is zero because
$n\geq1$. Finally $M<L$ since $L<1$.
\end{proof}

\subsection*{The determinant density}
For almost every $0<t<1$ define
\begin{align}
 A_0(t)&=L(1-t)^2\rho(t),\label{eq:A}\\
 N(t)&=\frac12\int_t^1(s-t/s)^2\rho(s)\rho(t/s)\,\frac{ds}{s},
 \qquad d(t)=A_0(t)-N(t).\label{eq:N}
\end{align}
Both $A_0$ and $N$ are integrable and nonnegative: by Tonelli,
$\int N=\frac12\int\!\int(s-u)^2\rho(s)\rho(u)\,ds\,du\leq M^2/2$.
Expanding \eqref{eq:Vmoment} and changing variables $t=su$ gives
\begin{equation}\label{eq:Dmoment}
 \frac{D_n}{256^n}=v_n^2-v_{n-1}v_{n+1}
                 =\int_0^1t^{n-1}d(t)\,dt\qquad(n\geq2).
\end{equation}
For clarity, if $a_k=\int u^k\rho(u)\,du$, the left side is
$L(a_{n-1}-2a_n+a_{n+1})-(a_{n-1}a_{n+1}-a_n^2)$;
the latter parenthesis is
$\frac12\int\!\int(su)^{n-1}(s-u)^2\rho(s)\rho(u)\,ds\,du$.
All changes of variables therefore involve finite nonnegative integrals.
The symmetry $s\leftrightarrow t/s$ also gives
\begin{equation}\label{eq:half}
 N(t)=\int_{\sqrt t}^1\frac{(a-t/a)^2}{a}\rho(a)\rho(t/a)\,da.
\end{equation}
We will prove $d(t)>0$ almost everywhere.

\subsection*{Positivity for $1/2\leq t<1$}
The series \eqref{eq:series} shows that $\rho$ decreases on $(1/2,1)$.
For $\sqrt t\leq a\leq1$, put $b=t/a$. Then $t\leq b\leq a$
and $\rho(b)\leq\rho(t)$. Moreover
$a-b\leq a(1-t)$, so $(a-b)^2/a\leq(1-t)^2$.
Consequently \eqref{eq:half} yields
\begin{equation}\label{eq:endbound}
 N(t)\leq M(1-t)^2\rho(t)<L(1-t)^2\rho(t)=A_0(t).
\end{equation}
The value at $t=1/2$ may be omitted.

\subsection*{Positivity for $0<t\leq1/4$}
For $0<u<1/2$, the descending Landen identity \cite[19.8.11--12]{dlmf}
in parameter convention is
\[
 K(4u(1-u))=\frac1{1-u}K\left(\frac{u^2}{(1-u)^2}\right).
\]
In DLMF's modulus convention use $k=2\sqrt{u(1-u)}$,
$k'=1-2u$, and $k_1=u/(1-u)$.
Differentiation and the series for $K$ give
\begin{equation}\label{eq:G}
 \rho(u)=\frac{u}{\pi(1-u)^2}
 G\left(\frac{u^2}{(1-u)^2}\right),\qquad
 G(z)=\sum_{j\geq0}\frac{2j+1}{j+1}c_jz^j.
\end{equation}
For example, the coefficient calculation uses
$(1+2j)c_j-2(j+1)c_{j+1}=(2j+1)c_j/(2(j+1))$.
Expanding the last factor in the following beta integral proves
\[
 G(z)=\frac2\pi\int_0^1
 \frac{r^{1/2}\,dr}{\sqrt{1-r}\sqrt{1-zr}}.
\]
It follows pointwise under the integral that
\begin{equation}\label{eq:Gratio}
 \frac{G(z_1)}{G(z_0)}\leq
 \sqrt{\frac{1-z_0}{1-z_1}}\qquad(0\leq z_0\leq z_1<1).
\end{equation}

We claim that whenever $0<b\leq a\leq1$ and $ab\leq1/4$,
\begin{equation}\label{eq:kernelrho}
 \frac{(a-b)^2}{a}\rho(b)\leq(1-ab)^2\rho(ab).
\end{equation}
Except at the irrelevant coincident midpoint, $b<1/2$.
By \eqref{eq:G}--\eqref{eq:Gratio}, the ratio of the left side
to the right side is at most
\[
 F(a,b)=\frac{(a-b)^2}{a^2(1-b)(1-ab)}
          \sqrt{\frac{1-2ab}{1-2b}}.
\]
For $a>b$, its logarithmic derivative in $a$ is
\[
 \partial_a\log F=\frac{2b}{a(a-b)}
       -\frac{ab^2}{(1-ab)(1-2ab)}>0.
\]
Indeed, its sign is that of $2-6t+5t^2-a^2t$, where $t=ab$,
and this is at least
$2-7t+5t^2=(1-t)(2-5t)>0$ on the required range.
If $b\leq1/4$, increasing $a$ to $1$ gives $F(a,b)\leq F(1,b)=1$.
If $1/4<b<1/2$, increase $a$ to $1/(4b)$ instead. There,
\[
 F(1/(4b),b)=\frac4{3\sqrt2}
     \frac{(1-2b)^{3/2}(1+2b)^2}{1-b}.
\]
Its logarithmic derivative has the sign of $2-15b+10b^2<0$
on $[1/4,1/2)$; at $b=1/4$ its value is $1$.
The case $a=b$ follows directly or by a limit. This proves
\eqref{eq:kernelrho}. Substituting $b=t/a$ in \eqref{eq:half}
and integrating against $\rho(a)\,da$ proves \eqref{eq:endbound}
also for $0<t\leq1/4$.

\subsection*{A finite certificate for $1/4\leq t<1/2$}
The Cauchy--Schwarz inequality for gamma integrals gives
$\Gamma(j+1/2)^2\leq\Gamma(j)\Gamma(j+1)$, hence
$c_j\leq1/(\pi j)$ for $j\geq1$.
With $m=4s(1-s)$, \eqref{eq:series} implies
\begin{equation}\label{eq:logbound}
 \rho(s)\leq\frac2{\pi^2}\sum_{j\geq1}\frac{m^j}{j+1}
 \leq-\frac2{\pi^2}\log(1-m)
 =\frac4{\pi^2}\bigl(-\log|1-2s|\bigr).
\end{equation}
In \eqref{eq:N}, use $2xy\leq x^2+y^2$ and the change
$s\leftrightarrow t/s$ in one of the two resulting integrals. This gives
\begin{equation}\label{eq:Ibound}
 N(t)\leq\frac8{\pi^4}\int_t^1 Q(s,t)\log^2|1-2s|\,ds,
 \qquad Q(s,t)=\frac{(s-t/s)^2}{s}.
\end{equation}
On the other hand \eqref{eq:G} and $G\geq1$ give
\begin{equation}\label{eq:Alower}
 A_0(t)\geq\frac{2t}{\pi^2}.
\end{equation}

For $j=4,5,6,7$, put $a_j=j/16$, and take
$t\in[a_j,(j+1)/16]$. For $i=4j,\ldots,63$, define the rational number
\[
 B_{j,i}=\max\{Q(s,t):s\in\{i/64,(i+1)/64\},\
                         t\in\{j/16,(j+1)/16\}\}.
\]
These four corner values bound $Q$ on the whole rectangle.
For fixed $s$, $Q$ is convex in $t$; for fixed $t$,
\[
 \partial_sQ=\frac{(s^2-t)(s^2+3t)}{s^4},
\]
so its only interior critical point in $s>0$ is a minimum.
Let
\[
 W_i=\int_{i/64}^{(i+1)/64}\log^2|1-2s|\,ds,\qquad
 I_j=\sum_{i=4j}^{63}B_{j,i}W_i.
\]
Enlarging the nonnegative integral in \eqref{eq:Ibound} down to $a_j$
shows that it is at most $I_j$. The exact certificate below proves
\begin{center}
\begin{tabular}{c|c|c}
 $j$ & certified rational upper bound for $I_j$ & $2a_j$\\\hline
 4 & $385126/1000000$ & $1/2$\\
 5 & $521335/1000000$ & $5/8$\\
 6 & $687221/1000000$ & $3/4$\\
 7 & $805114/1000000$ & $7/8$
\end{tabular}
\end{center}
Every middle entry is strictly below the corresponding right entry.
Since $\pi^2>8$, it follows that
\[
 N(t)<\frac{16a_j}{\pi^4}<\frac{2a_j}{\pi^2}
       \leq\frac{2t}{\pi^2}\leq A_0(t).
\]
Together with the endpoint arguments, this proves $d>0$ almost
everywhere on $(0,1)$.

Here are all details for verifying the finite table. Put
$\Phi(0)=0$ and
$\Phi(y)=y(\ell^2+2\ell+2)$ for $y>0$, where $\ell=-\log y$.
Thus $\Phi'(y)=\log^2y$.
Only arguments $y=k/32$, $0\leq k\leq32$, occur in the weights.
For $y>0$, set $z=(1-y)/(1+y)$ and
\begin{align*}
 \ell_-&=2\sum_{k=0}^{127}\frac{z^{2k+1}}{2k+1},\\
 \ell_+&=\ell_-+\frac{2z^{257}}{257(1-z^2)}.
\end{align*}
The power series $-\log y=2\operatorname{arctanh}z$ and a geometric
tail estimate prove $\ell_-\leq\ell\leq\ell_+$.
The function $y(\ell^2+2\ell+2)$ increases in $\ell\geq0$.
Round its lower bound down and upper bound up to multiples of $2^{-80}$,
obtaining $\Phi_-(y)$ and $\Phi_+(y)$. For $i<32$ use
\[
 W_i\leq\frac{\Phi_+((32-i)/32)-\Phi_-((31-i)/32)}2;
\]
for $i\geq32$ use
\[
 W_i\leq\frac{\Phi_+((i-31)/32)-\Phi_-((i-32)/32)}2.
\]
These are rational expressions with outward bounds throughout.
Multiplying by $B_{j,i}$, summing, and rounding up to the next
multiple of $10^{-6}$ gives the displayed table.
The accompanying standard-library Python program
\texttt{verify\_density\_certificate.py} implements these exact
instructions with \texttt{fractions.Fraction}; its complete rational
sums are retained in \texttt{density-certificate.json}.
This finite arithmetic certificate is a part of the proof.

\subsection*{Moment closure and all iterations}
\begin{proof}[Proof of Theorem \ref{thm:main}]
Define $\nu$ to be the image under $x=256t$ of the positive measure
with density $256^2t\,d(t)$ on $(0,1)$.
It is finite and has support $[0,256]$. Equation \eqref{eq:Dmoment}
gives $D_{m+2}=\int x^m\,d\nu(x)$ for every $m\geq0$.

For any finite positive measure $\eta$ supported on $[0,B]$, if
$a_m=\int x^m\,d\eta(x)$, then
\[
 a_ma_{m+2}-a_{m+1}^2
 =\frac12\int\!\int(xy)^m(x-y)^2\,d\eta(x)\,d\eta(y).
\]
This is again a moment sequence: use the image under multiplication
of $(x-y)^2\eta\otimes\eta/2$, supported on $[0,B^2]$.
Iteration proves infinite log-convexity.
If $\eta$ has full support $[0,B]$ with $B>0$, this new measure has
full support $[0,B^2]$: every point in $(0,B^2)$ is a product of two
distinct points of $(0,B)$, and small neighborhoods have positive
product measure and positive weight. Thus every iterate has positive
moments, proving the asserted strictness.
\end{proof}

\subsection*{Scope}
The proof covers all iteration orders and all required indices.
It uses the classical Landen identity, proved endpoint estimates,
and a reproducible exact finite certificate. It does not infer a
uniform inequality from numerical samples. The current published
source and live claim fields were checked before completion; the
separate literature audit specifies the search scope. No external
peer review or guaranteed novelty is claimed.

\section*{Task 2: Prove the optimal global density bound}
Determine the least constant $C$ for which
$N(t)\leq C(1-t)^2\rho(t)$ almost everywhere on $(0,1)$.
In particular, prove or refute the candidate $C=M=4/\pi-1$,
which the endpoint argument already establishes outside $(1/4,1/2)$.
Seek an analytic argument that also removes the finite certificate.

\section*{Task 3: Derive quantitative bounds for all iterates}
Use the explicit moment measures to bound
$(\mathcal L^r(D_{m+2}))_m$ above and below in terms of $m$ and $r$.
Obtain sharp asymptotics for fixed $r$ as $m\to\infty$, then determine
how far those estimates remain uniform when $r$ grows with $m$.

\section*{Task 4: Classify shifted determinant densities}
For each integer $h\geq2$, analyze
$V_n^2-V_{n-h}V_{n+h}$ for $n\geq h+1$.
Derive its analogue of $d(t)$ and determine which $h$ yield a
Stieltjes moment sequence after shifting the initial index.
Identify any failure by a rigorous obstruction or explicit counterexample.

\section*{Task 5: Extend the measure construction to beta laws}
Replace the arcsine variables in the preliminary lemma by independent
densities proportional to $(1-x^2)^{a-1}$ and $(1-y^2)^{b-1}$,
where $a,b>0$. With $U=(1+XY)/2$ and
$B_n=(2n+1)\E(U^n)-n\E(U^{n-1})$, determine the parameter region
where the determinant sequence $B_n^2-B_{n-1}B_{n+1}$ is a positive
moment sequence. Include the changed endpoint contributions.

\begin{thebibliography}{9}
\bibitem{source}
B. Y. Sun and J. J. Y. Zhao, \emph{Log-behavior of Two Sequences Related
to the Elliptic Integrals}, Acta Mathematicae Applicatae Sinica,
English Series \textbf{36} (2020), 590--602, Conjecture 5.1, p.~601.
\url{https://doi.org/10.1007/s10255-020-0952-1}.
Full publisher PDF:
\url{https://applmath.cjoe.ac.cn/jweb_yysxxb_en/EN/PDF/10.1007/s10255-020-0952-1}.
\bibitem{dlmf}
NIST Digital Library of Mathematical Functions, Section 19.8(ii),
descending Landen transformation, equations 19.8.11--12.
\url{https://dlmf.nist.gov/19.8#ii}.
\end{thebibliography}
\end{document}
